\begin{abstract}
Loss of the vertical support stiffness of a wheelset bearing, through loose fits, a degraded housing,
adapter or pad, or clearance growth, shifts the resonance of the axlebox on its support. On freight
wagons this resonance lies in the band of axle-box acceleration excited by rail roughness, so a
stiffness identified from it is confounded with the unknown roughness. We pool passes at
different speeds and payloads in a single debiased Whittle likelihood: the resonance stays fixed in
frequency while the roughness excitation moves with speed, so the roughness levels can be left as free
nuisance parameters. A finite-element wheelset is coupled to neighbouring
wheels through an infinite Timoshenko rail. For a representative freight wagon (Y25 bogies, ballasted
60E1 track) and 50 passes at 60--110~km/h, simulations give a bias of at most \HLBiasAbsMax\% and a
root-mean-square error of \HLRMSEMin--\HLRMSEMax\%, within \HLRatioCRBMax{} times the Cram\'er--Rao
bound. Ignoring wheel-to-wheel coupling biases the estimate by \IgnoreCouplingForty\%; nonlinear
Hertzian contact shifts it by only \ICDiffMean{} percentage points. Because absolute values carry a
model-form uncertainty of about \BudgetCoest\%, the method targets changes between periods. With the
foundation and contact stiffness re-estimated in each period, a decrease of
\MDCCalFiftyMin--\MDCCalFiftyMax\% is detectable from 50 passes per period and
\MDCCalTwoHundredMin--\MDCCalTwoHundredMax\% from 200, and a 20\% softer foundation, which appears as a
\DriftSoftTwentyApparent\% change if the track is held fixed, causes no bias. Clearance growth changes the
effective stiffness little (\EffPctTenEmpty\% per 10~\textmu m), so the method reports support-stiffness
loss rather than its mechanism.
\end{abstract}
